\documentclass[../main.tex]{subfiles}

\begin{document}
\section{Introduction}
\noindent  
% State the aim of your study and outline the importance of your study. You can consider both the importance of this dataset itself, but also the importance of comparing algorithms and their suitability for the task more generally.
\subsection{Aim}
This reports aims to compare three \gls{ML} algorithms on a medical image classification task involving the MedMNIST v2 dataset. While the overall collection includes 12 datasets of standardised medical images from carefully selected sources covering primary data modalities (e.g., X-ray, OCT, ultrasound, CT, electron microscope), the specific dataset used for this project is PathMNIST - a set of over 100,000 images from a study on colorectal cancer histology, utilising \gls{HE} stained images, classified into 9 different types of tissues\cite{yang2023medmnistv2}. 

\noindent Two given models were required to be analysed - both the \gls{MLP} and \gls{CNN}, while third model, the \gls{SVM}, was chosen by the team. Through completion of detailed analysis of the data, training the models, and comparing the results, an understanding of the strengths and weaknesses of each model in the given imaging task aimed to be developed.


\subsection{Importance}
Technology has long been at the forefront of medical advancements throughout history, and the use of \gls{ML} techniques in automating histopathological classification is just one example. Within itself, computational pathology using \gls{ML} enables the objective quantification of complex morphological features, significantly reducing the variability between observers and diagnostic subjectivity involved in manual examination.

Furthermore developing an understanding of how black box models compare to ones that are easier to understand like SVM's

This knowledge can be extrapolated across many other fields within medicine, whether it be disease and cancer identification, the discovery of novel predictive biomarkers, or the optimisation of drug screening workflows.
\todo{need to finish introduction}


\noindent 



\end{document}